% GENERATED FILE. Edit canonical lessons, then run npm run generate:books.
% canonical-source-hash: fnv1a64:2cd3260506f604c7
% canonical-lessons: TE-C181-goppa, TE-C181-vasana, TE-C181-asaktikaramaina, TE-C181-adbhutamaina, TE-C181-mamulu

\chapter{Great, Smell, Interesting, Wonderful, Ordinary}
\label{ch:te-great-smell-interesting-wonderful-ordinary}
\hlchaptermodality{181}

\begin{chapteropening}
\textbf{By the end of this chapter:} \emph{I can say great, a smell, interesting, wonderful and ordinary.}

Complete the last lesson of chapter 181: I can say great, a smell, interesting, wonderful and ordinary.
\end{chapteropening}

\section[\texorpdfstring{\te{గొప్ప} (goppa)}{గొప్ప (goppa)}]{\te{గొప్ప} (goppa) — great, eminent}

\label{lesson:TE-C181-goppa}



\begin{quote}
Before the new one: say the Telugu for pink, then the Telugu for orange.
\end{quote}

\subsection*{You'll want to know: \te{గొప్ప}}
\textbf{\te{గొప్ప}} — \emph{goppa} — \textquotedblleft{}great, eminent\textquotedblright{}.

\begin{grammarlens}[title={What you've built}]
One more word for this chapter.
\end{grammarlens}

\subsection*{Your turn}
\begin{itemize}
\raggedright
  \item \emph{Say it:} \emph{goppa}
  \item \emph{Say it:} \emph{goppa}, once more
  \item \emph{Say it:} say \emph{goppa}
  \item \emph{Recall:} say the Telugu for pink, then the Telugu for orange, then say \emph{goppa} again
\end{itemize}

\begin{tcolorbox}[breakable,colback=teal!4,colframe=teal!35!black,title={Before you move on}]
What does \te{గొప్ప} mean? (\textquotedblleft{}Great, eminent\textquotedblright{}.) Say it once more.
\end{tcolorbox}
\section[\texorpdfstring{\te{వాసన} (vāsana)}{వాసన (vāsana)}]{\te{వాసన} (vāsana) — a smell}

\label{lesson:TE-C181-vasana}



\begin{quote}
Before the new one: say the Telugu for orange, then the Telugu for great.
\end{quote}

\subsection*{You'll want to know: \te{వాసన}}
\textbf{\te{వాసన}} — \emph{vāsana} — \textquotedblleft{}a smell\textquotedblright{}.

\begin{grammarlens}[title={What you've built}]
One more word for this chapter.
\end{grammarlens}

\subsection*{Your turn}
\begin{itemize}
\raggedright
  \item \emph{Say it:} \emph{vāsana}
  \item \emph{Say it:} \emph{vāsana}, once more
  \item \emph{Say it:} say \emph{vāsana}
  \item \emph{Recall:} say the Telugu for orange, then the Telugu for great, then say \emph{vāsana} again
\end{itemize}

\begin{tcolorbox}[breakable,colback=teal!4,colframe=teal!35!black,title={Before you move on}]
What does \te{వాసన} mean? (\textquotedblleft{}A smell\textquotedblright{}.) Say it once more.
\end{tcolorbox}
\section[\texorpdfstring{\te{ఆసక్తికరమైన} (āsaktikaramaina)}{ఆసక్తికరమైన (āsaktikaramaina)}]{\te{ఆసక్తికరమైన} (āsaktikaramaina) — interesting}

\label{lesson:TE-C181-asaktikaramaina}



\begin{quote}
Before the new one: say the Telugu for great, then the Telugu for a smell.
\end{quote}

\subsection*{You'll want to know: \te{ఆసక్తికరమైన}}
\textbf{\te{ఆసక్తికరమైన}} — \emph{āsaktikaramaina} — \textquotedblleft{}interesting\textquotedblright{}.

\begin{grammarlens}[title={What you've built}]
One more word for this chapter.
\end{grammarlens}

\subsection*{Your turn}
\begin{itemize}
\raggedright
  \item \emph{Say it:} \emph{āsaktikaramaina}
  \item \emph{Say it:} \emph{āsaktikaramaina}, once more
  \item \emph{Say it:} say \emph{āsaktikaramaina}
  \item \emph{Recall:} say the Telugu for great, then the Telugu for a smell, then say \emph{āsaktikaramaina} again
\end{itemize}

\begin{tcolorbox}[breakable,colback=teal!4,colframe=teal!35!black,title={Before you move on}]
What does \te{ఆసక్తికరమైన} mean? (\textquotedblleft{}Interesting\textquotedblright{}.) Say it once more.
\end{tcolorbox}
\section[\texorpdfstring{\te{అద్భుతమైన} (adbhutamaina)}{అద్భుతమైన (adbhutamaina)}]{\te{అద్భుతమైన} (adbhutamaina) — wonderful, amazing}

\label{lesson:TE-C181-adbhutamaina}



\begin{quote}
Before the new one: say the Telugu for a smell, then the Telugu for interesting.
\end{quote}

\subsection*{You'll want to know: \te{అద్భుతమైన}}
\textbf{\te{అద్భుతమైన}} — \emph{adbhutamaina} — \textquotedblleft{}wonderful, amazing\textquotedblright{}.

\begin{grammarlens}[title={What you've built}]
One more word for this chapter.
\end{grammarlens}

\subsection*{Your turn}
\begin{itemize}
\raggedright
  \item \emph{Say it:} \emph{adbhutamaina}
  \item \emph{Say it:} \emph{adbhutamaina}, once more
  \item \emph{Say it:} say \emph{adbhutamaina}
  \item \emph{Recall:} say the Telugu for a smell, then the Telugu for interesting, then say \emph{adbhutamaina} again
\end{itemize}

\begin{tcolorbox}[breakable,colback=teal!4,colframe=teal!35!black,title={Before you move on}]
What does \te{అద్భుతమైన} mean? (\textquotedblleft{}Wonderful, amazing\textquotedblright{}.) Say it once more.
\end{tcolorbox}
\section[\texorpdfstring{\te{మామూలు} (māmūlu)}{మామూలు (māmūlu)}]{\te{మామూలు} (māmūlu) — ordinary, usual}

\label{lesson:TE-C181-mamulu}



\begin{quote}
Before the new one: say the Telugu for interesting, then the Telugu for wonderful.
\end{quote}

\subsection*{You'll want to know: \te{మామూలు}}
\textbf{\te{మామూలు}} — \emph{māmūlu} — \textquotedblleft{}ordinary, usual\textquotedblright{}.

\begin{grammarlens}[title={What you've built}]
One more word for this chapter.
\end{grammarlens}

\subsection*{Your turn}
\begin{itemize}
\raggedright
  \item \emph{Say it:} \emph{māmūlu}
  \item \emph{Say it:} \emph{māmūlu}, once more
  \item \emph{Say it:} say \emph{māmūlu}
  \item \emph{Recall:} say the Telugu for interesting, then the Telugu for wonderful, then say \emph{māmūlu} again
\end{itemize}

\begin{tcolorbox}[breakable,colback=teal!4,colframe=teal!35!black,title={Before you move on}]
What does \te{మామూలు} mean? (\textquotedblleft{}Ordinary, usual\textquotedblright{}.) Say it once more.
\end{tcolorbox}
